%----------------------------------------------------------------------------------------
% GRAFIEKEN VAN GONIOMETRISCHE FUNCTIES. Alle opgaven zijn eigen werk.
%
% Opbouw: Start, leerstof (theorie, werkwijze of uitleg, voorbeeld, korte oefening), en op een
% nieuwe pagina de Opdrachten (door elkaar, stijgende moeilijkheid), Koppel en Test jezelf.
% Labels van dit hoofdstuk beginnen met h3-.
%----------------------------------------------------------------------------------------

\mychapter{3}{Goniometrische functies}

\lkr{Leerplan B+S: LPD 13-16 (goniometrische functies), het deel over de grafieken van $\sin$, $\cos$ en $\tan$, de algemene sinusfunctie en periodieke modellen. De transformaties komen uit Algemene functies; hier krijgen ze namen (amplitude, periode, evenwichtslijn, faseverschuiving). Zonder grafisch rekentoestel: laat leerlingen schetsen vanuit vijf punten per periode.}

Eb en vloed, de lengte van de dag, de stand van een slinger, de spanning op het stopcontact: ze komen telkens terug, met een vast ritme. Zulke periodieke verschijnselen beschrijf je met de functies sinus en cosinus. In dit hoofdstuk teken je hun grafieken en pas je ze aan aan een situatie.

We vertrekken van de eenheidscirkel uit hoofdstuk 1: elke eigenschap van de grafieken volgt uit het punt dat over die cirkel draait.

\begin{oefeningen}[Start]
\oef{h3-start-1}

\oef{h3-verken}
\end{oefeningen}

%========================================================================================
\section{Sinus, cosinus en tangens als functie}

Laat de hoek $x$ (in radialen) doorlopen, en zet $\sin x$ uit als functiewaarde. Terwijl het punt over de eenheidscirkel draait, gaat zijn hoogte op en neer tussen $-1$ en $1$. Elke volle toer begint alles opnieuw.

\figuren{%
\begin{tikzpicture}[x=.62cm,y=.8cm,font=\scriptsize]
  \draw[black!14,very thin] (-3.14,-1) grid[xstep=1.5708,ystep=.5] (9.42,1);
  \draw[->,grijs] (-3.6,0) -- (9.9,0) node[right,text=black] {$x$};
  \draw[->,grijs] (0,-1.4) -- (0,1.5) node[above,text=black] {$y$};
  \foreach \t/\n in {-3.1416/-\pi,1.5708/\frac{\pi}{2},3.1416/\pi,4.7124/\frac{3\pi}{2},6.2832/2\pi,9.4248/3\pi}{\node[below,inner sep=1.5pt] at (\t,0) {$\n$};}
  \node[left] at (0,1) {$1$}; \node[left] at (0,-1) {$-1$};
  \draw[accent,thick] plot[domain=-3.1416:9.4248,samples=150,smooth] (\x,{sin(\x r)});
  \draw[accent,thick,dashed] plot[domain=-3.1416:9.4248,samples=150,smooth] (\x,{cos(\x r)});
  \node[accent,above] at (1.5708,1) {$\sin x$};
  \node[accent,above] at (6.2832,1) {$\cos x$};
\end{tikzpicture}}

\subsection{Periodieke functies}

Een grafiek die zich telkens na eenzelfde stuk herhaalt, noemen we periodiek. Het kleinste stuk waarna ze zich herhaalt, is de periode.

\begin{theorie}
Een functie $f$ is \textbf{\textcolor{accent}{periodiek}} met \textbf{\textcolor{accent}{periode}} $p > 0$ als $\formule{f(x + p) = f(x)}$ voor alle $x$ in het domein, en $p$ het kleinste positieve getal is waarvoor dat geldt.
\end{theorie}

Waarom het kleinste getal? Herhaalt een grafiek zich na $p$, dan ook na $2p$, $3p$, \ldots: ga twee keer $p$ verder, en je komt twee keer op dezelfde waarde uit. De periode is de kortste herhaling. Ken je de waarden op één periode, dan ken je de hele functie.

\voorbeeld{}{Een functie $f$ is periodiek met periode $3$, en $f(1) = 2$. Bereken $f(7)$ en $f(-5)$.\label{vb:h3-periodiek}}{%
\uitleg{We springen telkens een periode verder of terug, tot we bij $x = 1$ uitkomen, waar we de waarde kennen.}
\begin{stappenbreed}
f(7) = f(1 + 2 \cdot 3) = f(1) = 2 & Twee periodes verder: dezelfde waarde. \\[1ex]
f(-5) = f(1 - 2 \cdot 3) = f(1) = 2 & Twee periodes terug: ook dezelfde waarde.
\end{stappenbreed}}

\begin{oefeningen}
\oef{h3-periodiek-k}
\end{oefeningen}

\needspace{14\baselineskip}
\subsection{De sinusfunctie}

De eigenschappen van de sinusfunctie lees je af op de eenheidscirkel: $\sin x$ is de $y$-coördinaat van het punt $P$ bij de hoek $x$.

\begin{theorie}
\begin{stelling}{De sinusfunctie}
\renewcommand{\arraystretch}{1.3}
\begin{tabular}{@{}l@{\hspace{6mm}}l@{}}
domein & $\formule{\R}$ \\
bereik & $\formule{[-1, 1]}$ \\
periode & $\formule{2\pi}$ \\
nulwaarden & $\formule{x = k\pi}$ \quad ($k \in \Z$)
\end{tabular}
\end{stelling}
\begin{bewijs}
Elk reëel getal is een hoek, dus het domein is $\R$. De $y$-coördinaat van een punt op de eenheidscirkel ligt tussen $-1$ en $1$, en elke waarde $a \in [-1, 1]$ komt voor: de rechte $y = a$ snijdt de cirkel.

Uit hoofdstuk 1 weten we $\sin(x + 2\pi) = \sin x$. Is er een kleiner getal $p > 0$ met $\sin(x + p) = \sin x$ voor alle $x$? Neem $x = \frac{\pi}{2}$:
\bewijsregel{\sin\bigl(\tfrac{\pi}{2} + p\bigr) = \sin\tfrac{\pi}{2} = 1 .}
Het enige punt van de cirkel met $y = 1$ is $(0, 1)$, dus $\frac{\pi}{2} + p = \frac{\pi}{2} + k \cdot 2\pi$, en $p = k \cdot 2\pi$ met $k \geqslant 1$. De periode is dus $2\pi$.

Ten slotte is $\sin x = 0$ als $P$ op de $x$-as ligt: in $(1, 0)$ of $(-1, 0)$, dus bij $x = k\pi$.
\end{bewijs}
\end{theorie}

Op één periode, van $0$ tot $2\pi$, stijgt de sinus van $0$ naar de top $1$ bij $\frac{\pi}{2}$, daalt naar het dal $-1$ bij $\frac{3\pi}{2}$, en stijgt terug naar $0$. Daarna herhaalt alles zich. Een horizontale rechte $y = a$ met $-1 < a < 1$ snijdt de grafiek dus twee keer per periode.

\voorbeeld{}{Hoeveel oplossingen heeft $\sin x = 0{,}4$ in $[0, 4\pi]$?\label{vb:h3-aflezen}}{%
\naast{%
\uitleg{We tekenen de sinus en de rechte $y = 0{,}4$, en tellen de snijpunten per periode.}
\begin{stappenvrij}{3.4cm}{4.2cm}
2 \text{ per periode} & $0{,}4$ ligt tussen $-1$ en $1$: de rechte snijdt elke golf twee keer. \\[1ex]
2 \text{ periodes} & $[0, 4\pi]$ is twee keer $2\pi$ lang. \\[1ex]
2 \cdot 2 = 4 & Vier oplossingen.
\end{stappenvrij}}{%
\begin{tikzpicture}[x=.36cm,y=.9cm,font=\scriptsize]
  \draw[->,grijs] (-.5,0) -- (13.2,0) node[right,text=black] {$x$};
  \draw[->,grijs] (0,-1.3) -- (0,1.4) node[above,text=black] {$y$};
  \draw[accent,thick] plot[domain=0:12.566,samples=120,smooth] (\x,{sin(\x r)});
  \draw[grijs] (0,.4) -- (12.8,.4) node[right,text=black] {$0{,}4$};
  \foreach \p in {.4115,2.7301,6.6947,9.0133}{\fill[accent] (\p,.4) circle (2pt);}
  \node[below] at (6.2832,0) {$2\pi$}; \node[below] at (12.566,0) {$4\pi$};
  \node[left] at (0,1) {$1$}; \node[left] at (0,-1) {$-1$};
\end{tikzpicture}}}

\begin{oefeningen}
\oef{h3-sin-k}
\end{oefeningen}

\needspace{14\baselineskip}
\subsection{De cosinusfunctie}

In de figuur bovenaan lijkt de grafiek van de cosinus (streepjes) precies op die van de sinus, maar $\frac{\pi}{2}$ naar links geschoven. Dat is geen toeval.

\begin{theorie}
\begin{stelling}{De cosinusfunctie}
Voor elke $x$ is \quad $\formule{\cos x = \sin\bigl(x + \tfrac{\pi}{2}\bigr)}$, \quad en dus

\vspace{1mm}
\renewcommand{\arraystretch}{1.3}
\begin{tabular}{@{}l@{\hspace{6mm}}l@{}}
domein & $\formule{\R}$ \\
bereik & $\formule{[-1, 1]}$ \\
periode & $\formule{2\pi}$ \\
nulwaarden & $\formule{x = \tfrac{\pi}{2} + k\pi}$ \quad ($k \in \Z$)
\end{tabular}
\end{stelling}
\begin{bewijs}
Met supplementaire en daarna complementaire hoeken is
\bewijsregel{\sin\bigl(x + \tfrac{\pi}{2}\bigr) = \sin\bigl(\pi - (x + \tfrac{\pi}{2})\bigr) = \sin\bigl(\tfrac{\pi}{2} - x\bigr) = \cos x .}
De grafiek van de cosinus is dus die van de sinus, over $\frac{\pi}{2}$ naar links verschoven. Een verschuiving verandert het domein, het bereik en de periode niet. De nulwaarden schuiven mee: $x + \frac{\pi}{2} = k\pi$ geeft $x = -\frac{\pi}{2} + k\pi$, en dat zijn dezelfde getallen als $\frac{\pi}{2} + k\pi$.
\end{bewijs}
\end{theorie}

De cosinus begint in een top: $\cos 0 = 1$. Ze daalt tot het dal $-1$ bij $\pi$, en stijgt terug tot $1$ bij $2\pi$.

\voorbeeld{}{Geef alle $x$ in $[-2\pi, 2\pi]$ met $\cos x = 0$, en alle $x$ in dat interval met $\cos x = -1$.\label{vb:h3-cos}}{%
\naast{%
\uitleg{De nulwaarden staan in het kader; we vullen gehele $k$ in tot we buiten het interval vallen. Het dal ligt bij $\pi$, en daarna om de periode $2\pi$.}
\begin{stappenvrij}{3.9cm}{3.7cm}
\tfrac{\pi}{2} + k\pi \colon\ -\tfrac{3\pi}{2},\ -\tfrac{\pi}{2},\ \tfrac{\pi}{2},\ \tfrac{3\pi}{2} & $k = -2, -1, 0, 1$. \\[1ex]
\pi + k \cdot 2\pi \colon\ -\pi,\ \pi & $k = -1, 0$.
\end{stappenvrij}}{%
\begin{tikzpicture}[x=.34cm,y=.8cm,font=\scriptsize]
  \draw[->,grijs] (-6.8,0) -- (7.1,0) node[right,text=black] {$x$};
  \draw[->,grijs] (0,-1.3) -- (0,1.4) node[above,text=black] {$y$};
  \draw[accent,thick] plot[domain=-6.2832:6.2832,samples=120,smooth] (\x,{cos(\x r)});
  \foreach \p in {-4.7124,-1.5708,1.5708,4.7124}{\fill[accent] (\p,0) circle (2pt);}
  \foreach \p in {-3.1416,3.1416}{\filldraw[draw=accent,fill=white,thick] (\p,-1) circle (2pt);}
  \node[below] at (-6.2832,0) {$-2\pi$}; \node[below] at (6.2832,0) {$2\pi$};
  \node[above] at (3.1416,0) {$\pi$};
\end{tikzpicture}}}

\begin{oefeningen}
\oef{h3-cos-k}
\end{oefeningen}

\needspace{14\baselineskip}
\subsection{De tangensfunctie}

De tangens is het quotiënt $\frac{\sin x}{\cos x}$, en dus de richtingscoëfficiënt van de rechte $OP$. Daar volgen haar eigenschappen uit.

\begin{theorie}
\begin{stelling}{De tangensfunctie}
\renewcommand{\arraystretch}{1.3}
\begin{tabular}{@{}l@{\hspace{6mm}}l@{}}
domein & $\formule{\R \setminus \bigl\{\tfrac{\pi}{2} + k\pi \mid k \in \Z\bigr\}}$ \\
bereik & $\formule{\R}$ \\
periode & $\formule{\pi}$ \\
nulwaarden & $\formule{x = k\pi}$ \quad ($k \in \Z$)
\end{tabular}
\end{stelling}
\begin{bewijs}
De tangens bestaat als $\cos x \neq 0$, dus behalve in de nulwaarden $\frac{\pi}{2} + k\pi$ van de cosinus. Ze is $0$ als $\sin x = 0$, dus in $x = k\pi$.

Elke waarde $a \in \R$ komt voor: de rechte $y = ax$ door $O$ snijdt de eenheidscirkel in een punt $P$ rechts van de $y$-as, en de hoek $x$ van dat punt heeft $\tan x = a$.

Uit hoofdstuk 1 weten we $\tan(x + \pi) = \tan x$. Is $p > 0$ zo dat $\tan(x + p) = \tan x$ voor alle $x$, neem dan $x = 0$: $\tan p = 0$, dus $p = k\pi$ met $k \geqslant 1$. De periode is dus $\pi$.
\end{bewijs}
\end{theorie}

\naast{%
Waar de cosinus $0$ wordt, staat de rechte $OP$ verticaal. Nadert $x$ van links naar $\frac{\pi}{2}$, dan wordt $OP$ steeds steiler en $\tan x$ steeds groter. Nadert $x$ van rechts, dan is $\cos x < 0$ en wordt $\tan x$ steeds sterker negatief. De rechten $x = \frac{\pi}{2} + k\pi$ zijn daarom \textbf{verticale asymptoten}. Tussen twee asymptoten stijgt de tangens van $-\infty$ naar $+\infty$, en zo neemt ze op elke tak elke waarde precies één keer aan.}{%
\begin{tikzpicture}[x=.5cm,y=.42cm,font=\scriptsize]
  \draw[->,grijs] (-5.2,0) -- (5.4,0) node[right,text=black] {$x$};
  \draw[->,grijs] (0,-3.3) -- (0,3.5) node[above,text=black] {$y$};
  \foreach \lo/\hi in {-4.4/-1.88,-1.27/1.27,1.88/4.4}{\draw[accent,thick] plot[domain=\lo:\hi,samples=40,smooth] (\x,{tan(\x r)});}
  \foreach \k in {-1,0,1}{\draw[grijs,dashed] ({\k*3.1416+1.5708},-3.3) -- ({\k*3.1416+1.5708},3.3);}
  \node[below right,inner sep=1pt] at (1.5708,0) {$\frac{\pi}{2}$};
  \node[below left,inner sep=1pt] at (-1.5708,0) {$-\frac{\pi}{2}$};
  \node[below,inner sep=2pt] at (3.1416,0) {$\pi$};
\end{tikzpicture}}

\voorbeeld{}{Hoeveel oplossingen heeft $\tan x = 2$ in $[0, 4\pi]$?\label{vb:h3-tan}}{%
\naast{%
\uitleg{Snijd de grafiek met de rechte $y = 2$. De tangens neemt elke waarde één keer aan per tak, en er is één tak per periode $\pi$.}
\begin{stappenvrij}{3.4cm}{4.2cm}
1 \text{ per periode} & Elke waarde één keer per tak. \\[1ex]
4 \text{ periodes} & $[0, 4\pi]$ is vier keer $\pi$ lang. \\[1ex]
4 & Vier oplossingen.
\end{stappenvrij}}{%
\begin{tikzpicture}[x=.36cm,y=.5cm,font=\scriptsize]
  \draw[->,grijs] (-.5,0) -- (13.2,0) node[right,text=black] {$x$};
  \draw[->,grijs] (0,-3.3) -- (0,3.4) node[above,text=black] {$y$};
  \foreach \lo/\hi in {0/1.27,1.88/4.41,5.01/7.55,8.15/10.69,11.29/12.566}{\draw[accent,thick] plot[domain=\lo:\hi,samples=40,smooth] (\x,{tan(\x r)});}
  \foreach \k in {0,1,2,3}{\draw[grijs,dashed] ({\k*3.1416+1.5708},-3.3) -- ({\k*3.1416+1.5708},3.3);}
  \draw[grijs] (0,2) -- (12.6,2) node[right,text=black] {$2$};
  \foreach \k in {0,1,2,3}{\fill[accent] ({1.1071+\k*3.1416},2) circle (2pt);}
  \node[below] at (6.2832,0) {$2\pi$}; \node[below] at (12.566,0) {$4\pi$};
\end{tikzpicture}}}

\begin{oefeningen}
\oef{h3-aantal}
\end{oefeningen}

%========================================================================================
\section{De algemene sinusfunctie}

Een getij, een daglengte of een trillende veer schommelt zelden tussen $-1$ en $1$, en heeft zelden periode $2\pi$. We passen de sinus daarom aan met vier getallen $a$, $b$, $c$ en $d$, in de \textbf{algemene sinusfunctie}
\[f(x) = a\sin\bigl(b(x - c)\bigr) + d \qquad \text{met } a > 0 \text{ en } b > 0 .\]
Elk getal heeft een eigen effect op de grafiek. Dat zijn de transformaties die je al kent: verticaal en horizontaal uitrekken, en verschuiven. We bekijken ze één voor één, en geven ze een naam.

\begin{theorie}
Een periodieke functie met grootste waarde $M$ en kleinste waarde $m$ heeft
\[\formule{\text{evenwichtslijn } y = \frac{M + m}{2}} \qquad \formule{\text{amplitude } \frac{M - m}{2}} .\]
\end{theorie}

De evenwichtslijn ligt midden tussen de toppen en de dalen. De amplitude is de afstand van een top tot de evenwichtslijn.

\needspace{14\baselineskip}
\subsection{De amplitude}

\naast{%
Vermenigvuldig je de sinus met $a > 0$, dan wordt elke functiewaarde $a$ keer zo groot: de grafiek wordt verticaal uitgerekt met factor $a$. De nulwaarden blijven op hun plaats, de toppen gaan naar $a$ en de dalen naar $-a$. In de figuur is $a = 2$ (volle lijn) en $a = \frac12$ (streepjes), met $\sin x$ in het grijs.}{%
\begin{tikzpicture}[x=.55cm,y=.62cm,font=\scriptsize]
  \draw[black!14,very thin] (0,-2) grid[xstep=1.5708,ystep=1] (6.2832,2);
  \draw[->,grijs] (-.3,0) -- (6.8,0) node[right,text=black] {$x$};
  \draw[->,grijs] (0,-2.4) -- (0,2.5) node[above,text=black] {$y$};
  \foreach \t in {-2,-1,1,2}{\node[left,inner sep=1.5pt] at (0,\t) {$\t$};}
  \node[below,inner sep=1.5pt] at (3.1416,0) {$\pi$}; \node[below,inner sep=1.5pt] at (6.2832,0) {$2\pi$};
  \draw[grijs,thick] plot[domain=0:6.2832,samples=80,smooth] (\x,{sin(\x r)});
  \draw[accent,thick] plot[domain=0:6.2832,samples=80,smooth] (\x,{2*sin(\x r)});
  \draw[accent,thick,dashed] plot[domain=0:6.2832,samples=80,smooth] (\x,{.5*sin(\x r)});
\end{tikzpicture}}

\begin{theorie}
\begin{stelling}{Amplitude}
De functie $\formule{a\sin x}$ (met $a > 0$) heeft bereik $\formule{[-a, a]}$, evenwichtslijn $y = 0$ en \textbf{\textcolor{accent}{amplitude}} $\formule{a}$.
\end{stelling}
\begin{bewijs}
Uit $-1 \leqslant \sin x \leqslant 1$ volgt, omdat $a > 0$, dat $-a \leqslant a\sin x \leqslant a$. De waarden $a$ en $-a$ worden bereikt in $\frac{\pi}{2}$ en $\frac{3\pi}{2}$. Dus $M = a$ en $m = -a$, zodat de evenwichtslijn $y = 0$ is en de amplitude $\frac{a - (-a)}{2} = a$.
\end{bewijs}
\end{theorie}

Is $a$ negatief, dan wordt de grafiek ook gespiegeld om de $x$-as: de golf begint dan dalend. De amplitude is dan $|a|$. We schrijven een sinusfunctie daarom liefst met $a > 0$.

\voorbeeld{}{Schets $f(x) = 3\sin x$ over één periode, en geef het bereik.\label{vb:h3-amplitude}}{%
\naast{%
\uitleg{De nulwaarden zijn die van de sinus. Alleen de hoogte verandert: de top ligt op $3$, het dal op $-3$.}
\begin{stappenvrij}{3.4cm}{4.2cm}
(0, 0),\ (\pi, 0),\ (2\pi, 0) & Nulwaarden, zoals bij $\sin x$. \\[1ex]
\bigl(\tfrac{\pi}{2}, 3\bigr),\ \bigl(\tfrac{3\pi}{2}, -3\bigr) & Top en dal, $3$ keer zo hoog. \\[1ex]
\ber f = [-3, 3] & Amplitude $3$.
\end{stappenvrij}}{%
\begin{tikzpicture}[x=.55cm,y=.42cm,font=\scriptsize]
  \draw[black!14,very thin] (0,-3) grid[xstep=1.5708,ystep=1] (6.2832,3);
  \draw[->,grijs] (-.3,0) -- (6.8,0) node[right,text=black] {$x$};
  \draw[->,grijs] (0,-3.4) -- (0,3.6) node[above,text=black] {$y$};
  \foreach \t in {-3,3}{\node[left,inner sep=1.5pt] at (0,\t) {$\t$};}
  \node[below,inner sep=1.5pt] at (3.1416,0) {$\pi$}; \node[below,inner sep=1.5pt] at (6.2832,0) {$2\pi$};
  \draw[accent,thick] plot[domain=0:6.2832,samples=80,smooth] (\x,{3*sin(\x r)});
  \foreach \p in {(0,0),(1.5708,3),(3.1416,0),(4.7124,-3),(6.2832,0)}{\fill[accent] \p circle (2pt);}
\end{tikzpicture}}}

\begin{oefeningen}
\oef{h3-amplitude-k}
\end{oefeningen}

\needspace{14\baselineskip}
\subsection{De periode}

\naast{%
In $\sin 2x$ loopt de hoek $2x$ twee keer zo snel als $x$. Een volle golf is dus al klaar als $2x$ van $0$ tot $2\pi$ gaat, dus als $x$ van $0$ tot $\pi$ gaat: de grafiek wordt horizontaal samengedrukt. In de figuur staat $\sin 2x$ (volle lijn) naast $\sin x$ (grijs).}{%
\begin{tikzpicture}[x=.55cm,y=.62cm,font=\scriptsize]
  \draw[black!14,very thin] (0,-1) grid[xstep=1.5708,ystep=1] (6.2832,1);
  \draw[->,grijs] (-.3,0) -- (6.8,0) node[right,text=black] {$x$};
  \draw[->,grijs] (0,-1.4) -- (0,1.6) node[above,text=black] {$y$};
  \node[below,inner sep=1.5pt] at (3.1416,0) {$\pi$}; \node[below,inner sep=1.5pt] at (6.2832,0) {$2\pi$};
  \draw[grijs,thick] plot[domain=0:6.2832,samples=80,smooth] (\x,{sin(\x r)});
  \draw[accent,thick] plot[domain=0:6.2832,samples=120,smooth] (\x,{sin(2*\x r)});
\end{tikzpicture}}

\begin{theorie}
\begin{stelling}{Periode}
De functie $\formule{\sin bx}$ (met $b > 0$) heeft periode $\formule{\frac{2\pi}{b}}$.
\end{stelling}
\begin{bewijs}
Eerst herhaalt de functie zich na $\frac{2\pi}{b}$:
\bewijsregel{\sin\Bigl(b\bigl(x + \tfrac{2\pi}{b}\bigr)\Bigr) = \sin(bx + 2\pi) = \sin bx .}
Is $p > 0$ zo dat $\sin\bigl(b(x + p)\bigr) = \sin bx$ voor alle $x$, stel dan $u = bx$. Omdat $x$ alle reële getallen doorloopt, doet $u$ dat ook, en $\sin(u + bp) = \sin u$ voor alle $u$. Het getal $bp$ is dus een veelvoud $k \cdot 2\pi$ met $k \geqslant 1$, zoals in het bewijs bij de sinusfunctie, en $p \geqslant \frac{2\pi}{b}$.
\end{bewijs}
\end{theorie}

Het getal $b$ zegt hoeveel volle golven er in een stuk van lengte $2\pi$ passen. Hoe groter $b$, hoe korter de periode. Is $b < 1$, dan wordt de golf langer.

\voorbeeld{}{Geef de periode van $f(x) = \sin\frac{\pi x}{4}$ en schets één periode.\label{vb:h3-periode}}{%
\naast{%
\uitleg{Schrijf de functie als $\sin bx$: hier is $b = \frac{\pi}{4}$. Verdeel de periode daarna in vier gelijke stukken: nul, top, nul, dal, nul.}
\begin{stappenvrij}{3.6cm}{4cm}
\frac{2\pi}{b} = 2\pi \cdot \frac{4}{\pi} = 8 & Delen door $\frac{\pi}{4}$ is vermenigvuldigen met $\frac{4}{\pi}$. \\[2ex]
(0, 0),\ (2, 1),\ (4, 0), & Om de $\frac84 = 2$ een volgend punt. \\
(6, -1),\ (8, 0) &
\end{stappenvrij}}{%
\begin{tikzpicture}[x=.48cm,y=.62cm,font=\scriptsize]
  \draw[black!14,very thin] (0,-1) grid[xstep=1,ystep=1] (8,1);
  \draw[->,grijs] (-.3,0) -- (8.6,0) node[right,text=black] {$x$};
  \draw[->,grijs] (0,-1.4) -- (0,1.6) node[above,text=black] {$y$};
  \foreach \t in {2,4,6,8}{\node[below,inner sep=1.5pt] at (\t,0) {$\t$};}
  \node[left,inner sep=1.5pt] at (0,1) {$1$}; \node[left,inner sep=1.5pt] at (0,-1) {$-1$};
  \draw[accent,thick] plot[domain=0:8,samples=80,smooth] (\x,{sin(pi*\x/4 r)});
  \foreach \p in {(0,0),(2,1),(4,0),(6,-1),(8,0)}{\fill[accent] \p circle (2pt);}
\end{tikzpicture}}}

\begin{oefeningen}
\oef{h3-periode}
\end{oefeningen}

\needspace{16\baselineskip}
\subsection{De faseverschuiving}

Vervang je $x$ door $x - c$, dan schuift de grafiek over $c$ naar rechts. Bij een sinusfunctie heet die verschuiving de \textbf{faseverschuiving}: de golf begint later, of vroeger als $c < 0$.

\begin{theorie}
\begin{stelling}{Faseverschuiving}
De grafiek van $\formule{\sin\bigl(b(x - c)\bigr)}$ is die van $\sin bx$, over $\formule{c}$ naar rechts verschoven. Het stijgende nulpunt $(0, 0)$ gaat naar $\formule{(c, 0)}$.
\end{stelling}
\begin{bewijs}
Stel $g(x) = \sin bx$ en $f(x) = \sin\bigl(b(x - c)\bigr) = g(x - c)$. Ligt $(x_0, y_0)$ op de grafiek van $g$, dan is $f(x_0 + c) = g(x_0) = y_0$, dus $(x_0 + c, y_0)$ ligt op de grafiek van $f$. Omgekeerd komt elk punt van de grafiek van $f$ zo uit een punt van de grafiek van $g$.
\end{bewijs}
\end{theorie}

Let op: de verschuiving lees je pas af als $b$ buiten de haakjes staat. In $\sin\bigl(2x - \frac{\pi}{2}\bigr)$ is de verschuiving niet $\frac{\pi}{2}$.

\voorbeeld{}{Over hoeveel is de grafiek van $f(x) = \sin\bigl(2x - \frac{\pi}{2}\bigr)$ verschoven ten opzichte van die van $\sin 2x$?\label{vb:h3-fase}}{%
\naast{%
\uitleg{We zetten $b = 2$ buiten de haakjes, zodat de vorm $\sin\bigl(b(x - c)\bigr)$ verschijnt.}
\begin{stappenvrij}{3.6cm}{4cm}
2x - \tfrac{\pi}{2} = 2\bigl(x - \tfrac{\pi}{4}\bigr) & $\frac{\pi}{2}$ gedeeld door $2$. \\[1ex]
c = \tfrac{\pi}{4} & Een verschuiving over $\frac{\pi}{4}$ naar rechts, niet over $\frac{\pi}{2}$. \\[1ex]
\bigl(\tfrac{\pi}{4}, 0\bigr) & Het stijgende nulpunt.
\end{stappenvrij}
Controle: $f\bigl(\frac{\pi}{4}\bigr) = \sin 0 = 0$, en de functie stijgt daar.}{%
\begin{tikzpicture}[x=.62cm,y=.62cm,font=\scriptsize]
  \draw[black!14,very thin] (0,-1) grid[xstep=.7854,ystep=1] (4.712,1);
  \draw[->,grijs] (-.3,0) -- (5.2,0) node[right,text=black] {$x$};
  \draw[->,grijs] (0,-1.4) -- (0,1.6) node[above,text=black] {$y$};
  \node[below,inner sep=1.5pt] at (.7854,0) {$\frac{\pi}{4}$};
  \node[below,inner sep=1.5pt] at (3.1416,0) {$\pi$};
  \draw[grijs,thick] plot[domain=0:4.712,samples=80,smooth] (\x,{sin(2*\x r)});
  \draw[accent,thick] plot[domain=0:4.712,samples=80,smooth] (\x,{sin((2*\x-1.5708) r)});
  \draw[->,accent] (0,.15) -- (.7854,.15);
\end{tikzpicture}}}

\begin{oefeningen}
\oef{h3-fase-k}
\end{oefeningen}

\needspace{16\baselineskip}
\subsection{De evenwichtslijn}

Tel je $d$ op bij een functie, dan schuift haar grafiek over $d$ naar boven (naar onder als $d < 0$). De golf schommelt dan rond de rechte $y = d$.

\begin{theorie}
\begin{stelling}{Evenwichtslijn}
De functie $\formule{a\sin\bigl(b(x - c)\bigr) + d}$ (met $a, b > 0$) heeft grootste waarde $\formule{d + a}$, kleinste waarde $\formule{d - a}$, \textbf{\textcolor{accent}{evenwichtslijn}} $\formule{y = d}$ en amplitude $\formule{a}$.
\end{stelling}
\begin{bewijs}
De sinus van een getal ligt tussen $-1$ en $1$ en bereikt die waarden. Zoals bij de amplitude is dan $-a \leqslant a\sin\bigl(b(x - c)\bigr) \leqslant a$, en tel je $d$ op:
\bewijsregel{d - a \leqslant f(x) \leqslant d + a ,}
met beide grenzen bereikt. Dus is $M = d + a$ en $m = d - a$, zodat
\bewijsregel{\frac{M + m}{2} = d \qquad \text{en} \qquad \frac{M - m}{2} = a .}
\end{bewijs}
\end{theorie}

Omgekeerd vind je $a$ en $d$ uit de grootste en de kleinste waarde, zonder het voorschrift te kennen.

\voorbeeld{}{Een sinusfunctie schommelt tussen $-1$ en $7$. Geef de evenwichtslijn en de amplitude.\label{vb:h3-evenwicht}}{%
\naast{%
\uitleg{De evenwichtslijn ligt midden tussen de grootste en de kleinste waarde. De amplitude is de afstand van de top tot die lijn.}
\begin{stappenvrij}{3.8cm}{3.8cm}
d = \frac{7 + (-1)}{2} = 3 & Het midden: evenwichtslijn $y = 3$. \\[2ex]
a = \frac{7 - (-1)}{2} = 4 & De helft van het verschil.
\end{stappenvrij}
Controle: $3 + 4 = 7$ en $3 - 4 = -1$.}{%
\begin{tikzpicture}[x=.5cm,y=.3cm,font=\scriptsize]
  \draw[->,grijs] (-.3,0) -- (6.8,0) node[right,text=black] {$x$};
  \draw[->,grijs] (0,-1.6) -- (0,8) node[above,text=black] {$y$};
  \draw[grijs,dashed] (0,3) -- (6.4,3) node[right,text=black] {$y = 3$};
  \draw[lijngrijs] (0,7) -- (6.4,7); \draw[lijngrijs] (0,-1) -- (6.4,-1);
  \node[left,inner sep=1.5pt] at (0,7) {$7$}; \node[left,inner sep=1.5pt] at (0,-1) {$-1$};
  \draw[accent,thick] plot[domain=0:6.2832,samples=80,smooth] (\x,{4*sin(\x r)+3});
  \draw[<->,accent] (4.7124,3) -- node[right] {$4$} (4.7124,7);
\end{tikzpicture}}}

\begin{oefeningen}
\oef{h3-evenwicht-k}
\end{oefeningen}

\needspace{18\baselineskip}
\subsection{Een sinusgrafiek schetsen}

Nu zetten we de vier effecten samen. Begin bij het stijgende nulpunt $(0, 0)$ van $\sin x$. Door $b$ duurt een golf $\frac{2\pi}{b}$, door $c$ schuift dat punt naar $x = c$, door $a$ wordt de golf $a$ hoog, en door $d$ schuift alles naar de evenwichtslijn. Het punt $(0, 0)$ komt zo in $(c, d)$ terecht.

\begin{theorie}
\begin{stelling}{Algemene sinusfunctie}
De grafiek van $\formule{f(x) = a\sin\bigl(b(x - c)\bigr) + d}$ (met $a, b > 0$) is een sinusgrafiek met
\[\formule{\text{amplitude } a} \qquad \formule{\text{periode } \frac{2\pi}{b}} \qquad \formule{\text{evenwichtslijn } y = d} \qquad \formule{\text{beginpunt } (c, d)}.\]
In het \textbf{\textcolor{accent}{beginpunt}} snijdt de grafiek de evenwichtslijn stijgend.
\end{stelling}
\begin{bewijs}
Amplitude en evenwichtslijn staan in de vorige stelling. De periode is die van $\sin bx$, want verschuiven verandert de periode niet. Het beginpunt: $f(c) = a\sin 0 + d = d$, en net na $c$ is $b(x - c)$ een klein positief getal, waar de sinus stijgt.
\end{bewijs}
\end{theorie}

\voorbeeld{}{Schets $f(x) = 3\sin\bigl(2(x - \frac{\pi}{4})\bigr) + 1$ over één periode.\label{vb:h3-schets}}{%
\naast{%
\uitleg{Lees amplitude, periode, evenwichtslijn en beginpunt af. Verdeel de periode in vier gelijke stukken: evenwicht, top, evenwicht, dal, evenwicht.}
\begin{stappenvrij}{3.6cm}{3.8cm}
a = 3,\ \tfrac{2\pi}{b} = \pi,\ d = 1 & Amplitude, periode, evenwicht. \\
(c, d) = \bigl(\tfrac{\pi}{4}, 1\bigr) & Beginpunt. \\
\text{top } \bigl(\tfrac{\pi}{2}, 4\bigr) & Een kwart periode later. \\
\text{dal } \bigl(\pi, -2\bigr) & Driekwart periode. \\
\text{einde } \bigl(\tfrac{5\pi}{4}, 1\bigr) & Na een volle periode.
\end{stappenvrij}}{%
\begin{tikzpicture}[x=1.15cm,y=.42cm,font=\scriptsize]
  \draw[black!14,very thin] (0,-2) grid[xstep=.7854,ystep=1] (4.1,4);
  \draw[->,grijs] (-.2,0) -- (4.4,0) node[right,text=black] {$x$};
  \draw[->,grijs] (0,-2.4) -- (0,4.6) node[above,text=black] {$y$};
  \foreach \t/\n in {.7854/\frac{\pi}{4},1.5708/\frac{\pi}{2},3.1416/\pi,3.927/\frac{5\pi}{4}}{\node[below,inner sep=1.5pt] at (\t,0) {$\n$};}
  \foreach \t in {-2,1,4}{\node[left,inner sep=1.5pt] at (0,\t) {$\t$};}
  \draw[grijs,dashed] (0,1) -- (4.1,1);
  \draw[accent,thick] plot[domain=0:4.1,samples=100,smooth] (\x,{3*sin(2*(\x-.7854) r)+1});
  \foreach \p in {(.7854,1),(1.5708,4),(2.3562,1),(3.1416,-2),(3.927,1)}{\fill[accent] \p circle (1.6pt);}
\end{tikzpicture}}}

\begin{werkwijze}
Om een sinusgrafiek te schetsen: zet zo nodig $b$ buiten de haakjes. Teken de evenwichtslijn $y = d$, zet het beginpunt $(c, d)$, en daarna om de kwart periode een top ($d + a$), een evenwichtspunt, een dal ($d - a$) en een evenwichtspunt. Verbind met een golvende lijn en herhaal links en rechts.
\end{werkwijze}

\begin{oefeningen}
\oef{h3-kenmerken}

\oef{h3-schets}
\end{oefeningen}

\needspace{14\baselineskip}
\subsection{Een voorschrift uit de grafiek}

Omgekeerd haal je uit een grafiek het voorschrift: de evenwichtslijn ligt midden tussen top en dal, de amplitude is de helft van het verschil, en de periode is twee keer de afstand van een top tot het volgende dal.

Voor de faseverschuiving is er een handige keuze. Een sinus begint in een stijgend evenwichtspunt, een cosinus in een top. En omdat $\cos x = \sin\bigl(x + \frac{\pi}{2}\bigr)$, is ook $a\cos\bigl(b(x - c)\bigr) + d$ een sinusfunctie, met dezelfde amplitude, periode en evenwichtslijn. Ken je een top, dan is een cosinus met $c$ op die top het eenvoudigst.

\voorbeeld{}{Een sinusgrafiek heeft een top in $(1, 5)$ en het eerstvolgende dal in $(4, -1)$. Geef een voorschrift.\label{vb:h3-voorschrift}}{%
\naast{%
\uitleg{We kennen een top, dus we kiezen een cosinus met $c = 1$. De rest volgt uit de top en het dal.}
\begin{stappenvrij}{3.9cm}{3.8cm}
d = \tfrac{5 + (-1)}{2} = 2 & Evenwicht: midden tussen top en dal. \\[1ex]
a = \tfrac{5 - (-1)}{2} = 3 & Amplitude. \\[1ex]
\tfrac{2\pi}{b} = 2 \cdot (4 - 1) = 6 & Van top tot dal is een halve periode. \\[1ex]
b = \tfrac{2\pi}{6} = \tfrac{\pi}{3} & \\[1ex]
f(x) = 3\cos\bigl(\tfrac{\pi}{3}(x - 1)\bigr) + 2 & Top in $x = 1$.
\end{stappenvrij}}{%
\begin{tikzpicture}[x=.42cm,y=.32cm,font=\scriptsize]
  \draw[black!14,very thin] (-1,-1) grid (7,5);
  \draw[->,grijs] (-1.3,0) -- (7.5,0) node[right,text=black] {$x$};
  \draw[->,grijs] (0,-1.5) -- (0,5.8) node[above,text=black] {$y$};
  \foreach \t in {1,4,7}{\node[below,inner sep=1.5pt] at (\t,0) {$\t$};}
  \foreach \t in {-1,2,5}{\node[left,inner sep=1.5pt] at (0,\t) {$\t$};}
  \draw[grijs,dashed] (-1,2) -- (7,2);
  \draw[accent,thick] plot[domain=-1:7,samples=80,smooth] (\x,{3*cos(pi/3*(\x-1) r)+2});
  \fill[accent] (1,5) circle (2.5pt); \fill[accent] (4,-1) circle (2.5pt);
  \filldraw[draw=accent,fill=white,thick] (-.5,2) circle (2.5pt);
\end{tikzpicture}}
Wil je een sinus, neem dan het stijgende evenwichtspunt (het open bolletje): het ligt een kwart periode vóór de top, in $x = 1 - \frac64 = -\frac12$. Dat geeft ook $f(x) = 3\sin\bigl(\frac{\pi}{3}(x + \frac12)\bigr) + 2$.}

\begin{werkwijze}
Lees de grootste waarde $M$ en de kleinste waarde $m$ af. Dan is $d = \frac{M + m}{2}$ en $a = \frac{M - m}{2}$. De periode is twee keer de afstand van een top tot het volgende dal; daaruit volgt $b$. Kies voor $c$ een top (met een cosinus) of een stijgend evenwichtspunt (met een sinus). Controleer met een punt van de grafiek.
\end{werkwijze}

\begin{oefeningen}
\oef{h3-sinusfunctie}

\oef{h3-voorschrift-grafiek}
\end{oefeningen}

%========================================================================================
\section{Periodieke verschijnselen}

In een context ken je meestal de grootste en de kleinste waarde, en wanneer ze voorkomen. Daarmee stel je een sinusmodel op, net zoals een voorschrift uit een grafiek. Daarna bereken je met het model een waarde die je niet gemeten hebt.

\voorbeeld{}{In Brussel is de langste dag (21 juni, dag 172 van het jaar) ongeveer $16{,}4$ uur lang, de kortste (21 december, dag 355) ongeveer $8{,}1$ uur. Stel een model op voor de daglengte $L$ op dag $t$, en bereken de daglengte op 1 mei (dag 121).\label{vb:h3-model}}{%
\uitleg{De periode is een jaar. We kennen de dag van de top, dus we nemen een cosinus met de top op dag 172.}
\begin{stappenbreed}
d = \tfrac{16{,}4 + 8{,}1}{2} = 12{,}25,\quad a = \tfrac{16{,}4 - 8{,}1}{2} = 4{,}15 & Evenwicht en amplitude, in uur. \\[1ex]
b = \tfrac{2\pi}{365} & Periode 365 dagen. \\[1ex]
L(t) = 4{,}15\cos\bigl(\tfrac{2\pi}{365}(t - 172)\bigr) + 12{,}25 & Top op dag 172. \\[1ex]
L(121) = 4{,}15\cos\bigl(\tfrac{2\pi}{365} \cdot (-51)\bigr) + 12{,}25 \approx 14{,}9 & Ongeveer 14 uur 54 minuten.
\end{stappenbreed}
\figuren{\begin{tikzpicture}[x=.03cm,y=.22cm,font=\scriptsize]
  \draw[black!14,very thin] (0,8) grid[xstep=30.4167,ystep=2] (365,16);
  \draw[->,grijs] (0,8) -- (380,8) node[right,text=black] {$t$};
  \draw[->,grijs] (0,8) -- (0,17.5) node[above,text=black] {$L$};
  \foreach \t in {100,200,300}{\node[below,inner sep=1.5pt] at (\t,8) {$\t$};}
  \foreach \t in {8,12,16}{\node[left,inner sep=1.5pt] at (0,\t) {$\t$};}
  \draw[grijs,dashed] (0,12.25) -- (365,12.25);
  \draw[accent,thick] plot[domain=0:365,samples=100,smooth] (\x,{4.15*cos(2*pi/365*(\x-172) r)+12.25});
  \fill[accent] (172,16.4) circle (2.2pt); \fill[accent] (355,8.1) circle (2.2pt);
  \fill[accent] (121,14.9) circle (2.2pt) node[above left] {$(121;\ 14{,}9)$};
\end{tikzpicture}}
Het model is een benadering: de echte daglengte volgt niet precies een sinus, en wijkt enkele minuten af. Het dal valt in het model op dag $172 + 182{,}5 = 354{,}5$, een halve dag vóór 21 december.}

\begin{werkwijze}
Bepaal uit de context de grootste en de kleinste waarde, en de periode. Bereken $d$, $a$ en $b$. Kies een cosinus met $c$ op het tijdstip van de top, of een sinus met $c$ op een stijgend evenwichtspunt. Zeg in de context wat de eenheden zijn, en controleer het model met een gegeven waarde.
\end{werkwijze}

\begin{oefeningen}
\oef{h3-model}

\oef{h3-model-2}
\end{oefeningen}

%========================================================================================
\opdrachtenkop

\begin{oefeningen}
\needspace{6\baselineskip}
\oef{h3-o9}

\oef{h3-o2}

\oef{h3-o11}

\oef{h3-o1}

\needspace{12\baselineskip}
\oef{h3-o10}

\needspace{7\baselineskip}
\oef{h3-o12}

\needspace{10\baselineskip}
\oef{h3-o3}

\needspace{6\baselineskip}
\oef{h3-o13}

\needspace{9\baselineskip}
\oef{h3-o14}

\needspace{10\baselineskip}
\oef{h3-o16}

\oef{h3-o21}

\oef{h3-o4}

\needspace{9\baselineskip}
\oef{h3-o18}

\oef{h3-o17}

\oef{h3-o5}

\oef{h3-o20}

\oef{h3-o23}

\oef{h3-o19}

\oef{h3-o6}

\needspace{9\baselineskip}
\oef{h3-o24}

\oef{h3-o22}

\oef{h3-o7}

\needspace{12\baselineskip}
\oef{h3-o29}

\oef{h3-o26}

\oef{h3-o8}

\oef{h3-o30}

\oef{h3-o27}
\end{oefeningen}

%----------------------------------------------------------------------------------------
% KOPPEL: verbanden met andere leerstof
%----------------------------------------------------------------------------------------
\needspace{20\baselineskip}
\opdrachtsectie{Koppel}

\lkr{Afgeleiden en Integralen: kijk na wat al gezien is. Het pakket legt de volgorde niet vast.}

\begin{oefeningen}
\oef{h3-koppel-1}

\oef{h3-koppel-2}

\end{oefeningen}

%----------------------------------------------------------------------------------------
\opdrachtsectie{Test jezelf}

\begin{oefeningen}
\oef{h3-t1}

\oef{h3-t5}

\oef{h3-t2}

\oef{h3-t7}

\oef{h3-t3}

\oef{h3-t6}

\oef{h3-t4}
\end{oefeningen}

\herhaalschema{%
  \fv{oef:h3-t1}{vb:h3-aflezen}{oef:h3-aantal}%
  \fv{oef:h3-t5}{vb:h3-tan}{oef:h3-periode}%
  \fv{oef:h3-t2}{vb:h3-schets}{oef:h3-kenmerken}%
  \fv{oef:h3-t7}{vb:h3-schets}{oef:h3-schets}%
  \fv{oef:h3-t3}{vb:h3-voorschrift}{oef:h3-sinusfunctie}%
  \fv{oef:h3-t6}{vb:h3-voorschrift}{oef:h3-voorschrift-grafiek}%
  \fv{oef:h3-t4}{vb:h3-model}{oef:h3-model-2}}

\ifzelfstudie
\section*{Maak je eigen samenvatting}
Sluit dit boek en schrijf in je schrift: de grafieken van $\sin$, $\cos$ en $\tan$ met domein, bereik, periode en nulwaarden, de betekenis van $a$, $b$, $c$ en $d$ in $a\sin\bigl(b(x - c)\bigr) + d$, hoe je zo'n grafiek schetst met vijf punten, en hoe je uit een top en een dal een voorschrift opstelt. Vergelijk daarna met de kaders.
\fi
